\chapter{Comparative Synthesis}
\label{ch:synthesis}

Chapters~\ref{ch:representation} and~\ref{ch:supervision} analyzed representation and supervision 
separately, and Chapter~\ref{ch:Evaluation} gave a way to rate the results. This chapter discusses their combined impact on the final output.
The goal is to compare the methods directly, explain why some 
produce clean, editable SVGs while others produce visually 
rich outputs that fall apart when you try to edit them.
Three aspects cover the chapter: how representation and 
supervision interact to shape what comes out, where the 
realism-editability trade-off shows itself, and 
what conditions are required to produce usable SVG output.


\section{How Representation and Supervision Shape Structure}
\label{sec:rep_sup_interaction}

As already mentioned in the previous chapters, 
representation and supervision do not work independently. 
Choices of each one depend on the other and together they determine the structure of the output.  
Changing the representation under the same supervision signal 
produces a completely different output. Also, keeping the same representation and changing the supervision signal 
immediately impacts the result. 
Table~\ref{tab:representation_outcomes} maps these combinations across 
all analyzed models.


\begin{table}[htbp]
\centering
\caption{Representation type, supervision signal, 
and observed structural outcomes across the 
analyzed papers.}
\label{tab:representation_outcomes}
\footnotesize
\begin{tabularx}{\textwidth}{l l >{\raggedright\arraybackslash}p{2.7cm} >{\raggedright\arraybackslash}X}
\toprule
\textbf{Representation} & \textbf{Supervision} & 
\textbf{Models} & \textbf{Structural Outcome} \\
\midrule
Parametric    & Pixel recon.   & LIVE, SAMVG       
              & High path count, 
                fragmented structure \\
Parametric    & SDS            & VectorFusion      
              & Semantically rich, dense primitives, 
                not editable \\
Parametric    & VPSD           & SVGDreamer, SVGDreamer++ 
              & Improved diversity, object-aware 
                decomposition, dense \\
Neural implicit & SDS          & NIVeL, NeuralSVG  
              & Fixed layer count, constrained 
                structure, moderate editability \\
Sequential tokens & Reconstruction & DeepSVG, DeepIcon 
              & Clean paths within training domain \\
Neural path latent & VSD       & T2V-NPR           
              & Smooth, stable paths, 
                dataset-bounded \\
Continuous latent & Denoising  & SVGFusion         
              & Logical construction order, 
                layers, fast inference \\
Layered sequences & Denoising  & LayerTracer       
              & Designer-aligned layers, explicit 
                construction logic, strong editability \\
\bottomrule
\end{tabularx}
\end{table}

Diffusion-based supervision can be a very good example. VectorFusion 
and SVGDreamer both use it to optimize explicit B\'{e}zier 
paths~\cite{jainVectorFusionTexttoSVGAbstracting2022, 
xingSVGDreamerTextGuided2024}. They are run with large path budgets, 
up to hundreds of primitives, because that is how unconstrained 
parametric paths approximate fine visual detail.
However, NeuralSVG uses the same kind of signal running it through 
an implicit neural representation and gets comparable 
visual quality with 16 
shapes (Figure~\ref{fig:supervision_comparison}),
and fewer primitives or shapes mean a better SVG file.
The representation is not the only difference. NeuralSVG also adds a LoRA adapter, 
saliency-based initialization and dropout-based ordering, and its shape count is set 
in advance~\cite{polaczekNeuralSVGImplicitRepresentation2025}. But the representation is the 
biggest of these differences, and it changes the structural character of the output.

Reconstruction losses show the same example from a different 
angle. LIVE and SAMVG use the same kind of pixel loss, but SAMVG starts 
from segmentation masks and ends up with cleaner boundaries and fewer 
paths (Section~\ref{sec:init_priors}). 
That means where optimization started also made the difference.

Dataset-driven methods are constrained, so some approaches may have similar --- icon-wise and simple results. 
Representation is fixed by design and whatever structure 
shows up in the output is learned from the training data.
DeepSVG and DeepIcon produce clean, compact 
icons because the SVG-Icons8 dataset contains simple graphics 
with small path 
counts~\cite{carlierDeepSVGHierarchicalGenerative2020}. The 
reconstruction loss teaches them to be tidy through the data.
LayerTracer pushes this further: its training data 
encodes step-by-step construction order, so the model picks 
up layering as a convention rather than discovering it through 
optimization~\cite{songLayerTracerCognitiveAlignedLayered2025}. 
SVGFusion and T2V-NPR sit in between, combining learned 
latent structure with semantic supervision. Their outputs are 
more coherent than unconstrained optimization but still 
bounded by what the training corpus 
contains~\cite{xingSVGFusionScalableTexttoSVG2025, 
zhangTexttoVectorGenerationNeural2024}.

Representations that constrain 
decomposition (implicit fields, learned latent spaces, layered 
sequences) produce cleaner outputs than 
representations that do not limit primitive count, 
regardless of the supervision signal 
(Tables~\ref{tab:rep_structure} and~\ref{tab:representation_outcomes}). And supervision operating 
within a structured representation space produces more stable 
geometry than supervision operating directly in pixel space. 
NeuralSVG with SDS inside an implicit representation and 
T2V-NPR with VSD inside a learned path latent are both 
examples of this~\cite{polaczekNeuralSVGImplicitRepresentation2025, 
zhangTexttoVectorGenerationNeural2024}.

\section{The Realism--Editability Trade-off}
\label{sec:tradeoff}

The most consistent pattern across models is the 
tension between visual realism and editability. 
The tension is observable in specific models.
Methods driven by diffusion-based supervision achieve visually 
detailed outputs by refining geometry until the rendered image 
aligns with learned image priors. This refinement increases 
structural complexity, but undermines usability at the same time. VectorFusion 
uses 64 paths per SVG~\cite{jainVectorFusionTexttoSVGAbstracting2022}. LIVE shows the 
same problem without any diffusion model: it represents a complex image with up to 
256 paths added layer by layer~\cite{maLayerwiseImageVectorization2022}. In such a file, 
selecting or modifying a single visual component requires 
unlimited patience. SVGDreamer++ explicitly acknowledges this 
failure, noting that optimization-based approaches produce 
outputs where objects are entangled across paths, making 
component-level editing effectively 
impossible~\cite{xingSVGDreamerAdvancingEditability2025}.
However, this structural cost comes with a visual benefit, 
as Figure~\ref{fig:realism_editability} illustrates.



\begin{figure}[htbp]
    \centering
    \begin{tabular}{cc}
    \textbf{SVGDreamer} & \textbf{SVGFusion} \\
    \small{256--512 paths, VPSD} & 
    \small{few shapes, denoising} \\[4pt]
    \includegraphics[width=0.36\textwidth]{svgdreamer_example.png} &
    \includegraphics[width=0.36\textwidth]{svgfusion_example.png} \\
    \end{tabular}
    \caption{The realism-editability trade-off. SVGDreamer 
    (left) produces visually detailed output through 
    hundreds of optimized paths that overlap and cannot be 
    individually selected in a design tool. SVGFusion 
    (right) generates a layered SVG where each visual 
    component is a separate, editable element, at the cost 
    of reduced photorealistic 
    detail~\cite{xingSVGDreamerTextGuided2024,
    xingSVGFusionScalableTexttoSVG2025}.}
    \label{fig:realism_editability}
\end{figure}

SVGDreamer and SVGDreamer++ consistently produce the most 
visually rich and semantically detailed outputs among the 
analyzed methods~\cite{xingSVGDreamerTextGuided2024,xingSVGDreamerAdvancingEditability2025}.
Their large path budgets, guided by strong diffusion priors, 
fragment the structure, and that is precisely what enables 
photorealistic quality.

The correct choice of models and methods depends 
on whether the use prioritizes visual output quality 
or editability. Table~\ref{tab:framework} shows this trade-off 
for the parametric methods: VectorFusion and SVGDreamer are rated 
+ on semantic alignment but $-$ or \textasciitilde{} on path economy. 
The neural implicit methods are the exception.
Models that constrain representation tend to produce outputs where fewer primitives are used. 
NeuralSVG represents complex scenes using a small 
fixed number of shapes, improving usability despite reduced 
photorealistic detail~\cite{polaczekNeuralSVGImplicitRepresentation2025}. 
But even though it is constrained to a certain number of shapes, 
its visual quality, as shown in Figure~\ref{fig:supervision_comparison}, 
does not fall behind the compared models, in terms of photorealistic quality.


\section{Conditions for Clean and Editable SVG}
\label{sec:conditions}

The analysis in Sections~\ref{sec:rep_sup_interaction} 
and~\ref{sec:tradeoff} points to five conditions that 
consistently separate editable outputs from fragmented 
ones. Table~\ref{tab:conditions} lists them with the 
models that satisfy each.

\begin{table}[htbp]
\centering
\caption{Five conditions for clean and editable SVG 
output, derived from the representation and supervision 
analysis in Chapters~\ref{ch:representation} 
and~\ref{ch:supervision}.}
\label{tab:conditions}
\small
\begin{tabularx}{\textwidth}{>{\raggedright\arraybackslash}p{3.6cm} >{\raggedright\arraybackslash}X >{\raggedright\arraybackslash}p{4.2cm}}
\hline
\textbf{Condition} & \textbf{Models} & \textbf{Where analyzed} \\
\hline
Primitive count control & NeuralSVG, NIVeL, T2V-NPR, SVGDreamer, SVGDreamer++ & Section~\ref{sec:parametric}, Table~\ref{tab:rep_structure} \\
Semantic decomposition & SAMVG, NeuralSVG, LayerTracer & Section~\ref{sec:init_priors} \\
Constrained supervision & NeuralSVG, T2V-NPR & Section~\ref{sec:diffusion_supervision} \\
Informed initialization & SAMVG, SVGDreamer, NeuralSVG & Section~\ref{sec:init_priors} \\
Design-aware priors & LayerTracer, SVGFusion & Section~\ref{sec:dataset_vs_instance} \\
\hline
\end{tabularx}
\end{table}

No model satisfies all five. NeuralSVG comes closest 
on the optimization side. LayerTracer comes closest on 
the dataset side. NeuralSVG is more flexible but slower. 
LayerTracer is faster but limited to what its training 
data covers.

Domain coverage is a related limitation. Dataset-driven 
models only generate what they were trained on: icons 
(DeepSVG, DeepIcon) or 
the SVGX corpus 
(SVGFusion)~\cite{carlierDeepSVGHierarchicalGenerative2020, 
bingDeepIconHierarchicalNetwork2024, 
xingSVGFusionScalableTexttoSVG2025}. Optimization-based 
methods handle broader domains through diffusion priors 
but with the structural costs described in 
Section~\ref{sec:rep_sup_interaction}. None of the 
thirteen covers typography, technical illustration, or 
character design.




